% ---------------------------------------------------------
\chapter{Einleitung}\label{chap:introduction}
% ---------------------------------------------------------

Forler, Lucks und Wenzel haben 2015 in \cite{PHC:Catena:2015} Catena
vorgestellt. Dabei handelt es sich um ein Verfahren zum hashen von
Passwörtern, welches sich nicht effizient auf einer \ac{GPU}
parallelisieren lässt. Wie eine moderne \ac{CPU} verfügen \acp{GPU}
über \ldots


% ---------------------------------------------------------
\section{Problemstellung und Motivation}
% ---------------------------------------------------------

\lipsum[1-2]

\paragraph{Quellen.} 
Bei der \ac{OMG} handelt es sich um einen Konsortium welches sich mit
der objektorientierten Programmierung beschäftigt.\cite{OMG2020}

In seiner Doktoarbeit \cite{DBLP:phd/dnb/Forler15} hat Herr Forler
\ldots


In \cite{knuth97} beschäftigt sich Donald E. Knuth mit Algorihmen und
Datenstrukturen.


In \cite{DBLP:conf/africacrypt/FleischmannFL12} haben Fleischmann,
Forler und Lucks einen Sicherheitsbeweis für MDC-4 vorgestellt.


% ---------------------------------------------------------
\section{Ziel der Arbeit}
% ---------------------------------------------------------



% ---------------------------------------------------------
\section{Aufbau der Arbeit}
% --------------------------------------------------------- 

In Kapitel \ref{chap:preliminaries} werden die Grunflagen vorgestellt
und Kapitel \ref{chap:concept} wird das Konzept der Anwendung
präsentiert. Wie diese Umgesetzt wurde wird in Kapitel
\ref{chap:implementation} erläutert. \ldots
